%%
%% This is file `acmtog.tex',
%% generated with the docstrip utility.
%%
%% The original source files were:
%%
%% samples.dtx  (with options: `all,acmtog')
%% 
%% IMPORTANT NOTICE:
%% 
%% For the copyright see the source file.
%% 
%% Any modified versions of this file must be renamed
%% with new filenames distinct from acmtog.tex.
%% 
%% For distribution of the original source see the terms
%% for copying and modification in the file samples.dtx.
%% 
%% This generated file may be distributed as long as the
%% original source files, as listed above, are part of the
%% same distribution. (The sources need not necessarily be
%% in the same archive or directory.)
%%
%%
%% Commands for TeXCount
%TC:macro \cite [option:text,text]
%TC:macro \citep [option:text,text]
%TC:macro \citet [option:text,text]
%TC:envir table 0 1
%TC:envir table* 0 1
%TC:envir tabular [ignore] word
%TC:envir displaymath 0 word
%TC:envir math 0 word
%TC:envir comment 0 0
%%
%% The first command in your LaTeX source must be the \documentclass
%% command.
%%
%% For submission and review of your manuscript please change the
%% command to \documentclass[manuscript, screen, review]{acmart}.
%%
%% When submitting camera ready or to TAPS, please change the command
%% to \documentclass[sigconf]{acmart} or whichever template is required
%% for your publication.
%%
%%
\documentclass[acmtog]{acmart}
%%
%% \BibTeX command to typeset BibTeX logo in the docs
\AtBeginDocument{%
  \providecommand\BibTeX{{%
    Bib\TeX}}}

%% Rights management information.  This information is sent to you
%% when you complete the rights form.  These commands have SAMPLE
%% values in them; it is your responsibility as an author to replace
%% the commands and values with those provided to you when you
%% complete the rights form.
\setcopyright{acmlicensed}
\copyrightyear{2018}
\acmYear{2018}
\acmDOI{XXXXXXX.XXXXXXX}

%% These commands are for a JOURNAL article.
\acmJournal{TOG}
\acmVolume{37}
\acmNumber{4}
\acmArticle{111}
\acmMonth{8}

%%
%% Submission ID.
%% Use this when submitting an article to a sponsored event. You'll
%% receive a unique submission ID from the organizers
%% of the event, and this ID should be used as the parameter to this command.
%%\acmSubmissionID{123-A56-BU3}

%%
%% For managing citations, it is recommended to use bibliography
%% files in BibTeX format.
%%
%% You can then either use BibTeX with the ACM-Reference-Format style,
%% or BibLaTeX with the acmnumeric or acmauthoryear sytles, that include
%% support for advanced citation of software artefact from the
%% biblatex-software package, also separately available on CTAN.
%%
%% Look at the sample-*-biblatex.tex files for templates showcasing
%% the biblatex styles.
%%
%%
%% The majority of ACM publications use numbered citations and
%% references.  The command \citestyle{authoryear} switches to the
%% "author year" style.
%%
%% If you are preparing content for an event
%% sponsored by ACM SIGGRAPH, you must use the "author year" style of
%% citations and references.
\citestyle{acmauthoryear}


%%
%% end of the preamble, start of the body of the document source.
\begin{document}

%%
%% The "title" command has an optional parameter,
%% allowing the author to define a "short title" to be used in page headers.
\title{SakunaGeoGraPH}

%%
%% The "author" command and its associated commands are used to define
%% the authors and their affiliations.
%% Of note is the shared affiliation of the first two authors, and the
%% "authornote" and "authornotemark" commands
%% used to denote shared contribution to the research.
\author{sample}
\email{sample@up.edu.ph}
\authornotemark[1]
\affiliation{%
  \institution{University of the Philippines Diliman}
  \city{Quezon City}
}
%%
%% By default, the full list of authors will be used in the page
%% headers. Often, this list is too long, and will overlap
%% other information printed in the page headers. This command allows
%% the author to define a more concise list
%% of authors' names for this purpose.
\renewcommand{\shortauthors}{Trovato et al.}

%%
%% The abstract is a short summary of the work to be presented in the
%% article.
\begin{abstract}
\end{abstract}

%%
%% The code below is generated by the tool at http://dl.acm.org/ccs.cfm.
%% Please copy and paste the code instead of the example below.
%%
\begin{CCSXML}
<ccs2012>
 <concept>
  <concept_id>00000000.0000000.0000000</concept_id>
  <concept_desc>Do Not Use This Code, Generate the Correct Terms for Your Paper</concept_desc>
  <concept_significance>500</concept_significance>
 </concept>
 <concept>
  <concept_id>00000000.00000000.00000000</concept_id>
  <concept_desc>Do Not Use This Code, Generate the Correct Terms for Your Paper</concept_desc>
  <concept_significance>300</concept_significance>
 </concept>
 <concept>
  <concept_id>00000000.00000000.00000000</concept_id>
  <concept_desc>Do Not Use This Code, Generate the Correct Terms for Your Paper</concept_desc>
  <concept_significance>100</concept_significance>
 </concept>
 <concept>
  <concept_id>00000000.00000000.00000000</concept_id>
  <concept_desc>Do Not Use This Code, Generate the Correct Terms for Your Paper</concept_desc>
  <concept_significance>100</concept_significance>
 </concept>
</ccs2012>
\end{CCSXML}

\ccsdesc[500]{Do Not Use This Code~Generate the Correct Terms for Your Paper}
\ccsdesc[300]{Do Not Use This Code~Generate the Correct Terms for Your Paper}
\ccsdesc{Do Not Use This Code~Generate the Correct Terms for Your Paper}
\ccsdesc[100]{Do Not Use This Code~Generate the Correct Terms for Your Paper}

%%
%% Keywords. The author(s) should pick words that accurately describe
%% the work being presented. Separate the keywords with commas.
\keywords{knowledge graphs, GeoSPARQL, disaster data, geocoding, data integration, Philippines}

\received{December 7, 2026}
\received[revised]{December 7, 2026}
\received[accepted]{December 7, 2026}

%%
%% This command processes the author and affiliation and title
%% information and builds the first part of the formatted document.
\maketitle

\section{Abstract}

lorem ipsum  lorem ipsum  lorem ipsum  lorem ipsum  lorem ipsum  lorem ipsum  lorem ipsum  lorem ipsum  lorem ipsum  lorem ipsum  lorem ipsum  lorem ipsum  lorem ipsum  lorem ipsum  lorem ipsum  lorem ipsum  lorem ipsum  lorem ipsum  lorem ipsum  lorem ipsum  lorem ipsum  lorem ipsum  lorem ipsum  lorem ipsum  lorem ipsum  lorem ipsum  lorem ipsum  lorem ipsum  lorem ipsum  lorem ipsum  lorem ipsum  lorem ipsum  lorem ipsum  lorem ipsum  lorem ipsum  lorem ipsum  lorem ipsum  lorem ipsum  lorem ipsum  lorem ipsum  lorem ipsum  lorem ipsum  lorem ipsum  lorem ipsum  lorem ipsum  lorem ipsum  lorem ipsum  lorem ipsum  lorem ipsum  lorem ipsum  lorem ipsum

\section{Introduction}

lorem ipsum  lorem ipsum  lorem ipsum  lorem ipsum  lorem ipsum  lorem ipsum  lorem ipsum  lorem ipsum  lorem ipsum  lorem ipsum  lorem ipsum  lorem ipsum  lorem ipsum  lorem ipsum  lorem ipsum  lorem ipsum  lorem ipsum  lorem ipsum  lorem ipsum  lorem ipsum  lorem ipsum  lorem ipsum  lorem ipsum  lorem ipsum  lorem ipsum  lorem ipsum  lorem ipsum  lorem ipsum  lorem ipsum  lorem ipsum  lorem ipsum  lorem ipsum  lorem ipsum  lorem ipsum  lorem ipsum  lorem ipsum  lorem ipsum  lorem ipsum  lorem ipsum  lorem ipsum  lorem ipsum  lorem ipsum  lorem ipsum  lorem ipsum  lorem ipsum  lorem ipsum  lorem ipsum  lorem ipsum  lorem ipsum  lorem ipsum  lorem ipsum  


\section{Review of Related Literature}

lorem ipsum  lorem ipsum  lorem ipsum  lorem ipsum  lorem ipsum  lorem ipsum  lorem ipsum  lorem ipsum  lorem ipsum  lorem ipsum  lorem ipsum  lorem ipsum  lorem ipsum  lorem ipsum  lorem ipsum  lorem ipsum  lorem ipsum  lorem ipsum  lorem ipsum  lorem ipsum  lorem ipsum  lorem ipsum  lorem ipsum  lorem ipsum  lorem ipsum  lorem ipsum  lorem ipsum  lorem ipsum  lorem ipsum  lorem ipsum  lorem ipsum  lorem ipsum  lorem ipsum  lorem ipsum  lorem ipsum  lorem ipsum  lorem ipsum  lorem ipsum  lorem ipsum  lorem ipsum  lorem ipsum  lorem ipsum  lorem ipsum  lorem ipsum  lorem ipsum  lorem ipsum  lorem ipsum  lorem ipsum  lorem ipsum  lorem ipsum  lorem ipsum

\section{Research Objective}

lorem ipsum  lorem ipsum  lorem ipsum  lorem ipsum  lorem ipsum  lorem ipsum  lorem ipsum  lorem ipsum  lorem ipsum  lorem ipsum  lorem ipsum  lorem ipsum  lorem ipsum  lorem ipsum  lorem ipsum  lorem ipsum  lorem ipsum  lorem ipsum  lorem ipsum  lorem ipsum  lorem ipsum  lorem ipsum  lorem ipsum  lorem ipsum  lorem ipsum  lorem ipsum  lorem ipsum  lorem ipsum  lorem ipsum  lorem ipsum  lorem ipsum  lorem ipsum  lorem ipsum  lorem ipsum  lorem ipsum  lorem ipsum  lorem ipsum  lorem ipsum  lorem ipsum  lorem ipsum  lorem ipsum  lorem ipsum  lorem ipsum  lorem ipsum  lorem ipsum  lorem ipsum  lorem ipsum  lorem ipsum  lorem ipsum  lorem ipsum  lorem ipsum  

\subsection{Research Objective 1}


\subsection{Research Objective 2}

\section{Data and Methodology}
\label{sec:method}

% Drafted from paper/paper_writing.md (sections 5.1-5.6), 25 Nov 2026.

\subsection{Overview}

This work extends SakunaGraPH, a knowledge graph that integrates Philippine disaster records
from heterogeneous sources into a single RDF graph, with two contributions: (1) a point
coordinate for every administrative unit in the Philippine Standard Geographic Code (PSGC)
down to the city and municipality level, and (2) a representative point (\emph{event
centroid}) for each disaster event, computed from the locations the event affected. Because
events in SakunaGraPH reference locations through PSGC identifiers rather than free text, the
coordinates are attached once, to the administrative units, and every event from every source
inherits them through its location links. The situation reports of the Department of Social
Welfare and Development's DROMIC~\cite{dswd2026dromic} are the primary event source studied,
since they are the most frequently updated.

The method has three stages (Figure~\ref{fig:overview}): (1) build an RDF graph of all PSGC
units from the PSA publication datafile; (2) derive one point per region, province, city,
municipality and sub-municipality from official boundary polygons and join it to the PSGC
graph; (3) for each event, aggregate the points of its affected municipalities into one event
point, weighted by the number of affected people.

\begin{figure}[t]
  \centering
  \fbox{\parbox[c][3.2cm][c]{0.9\columnwidth}{\centering\small
    [Figure F1 placeholder] PSGC backbone $\rightarrow$ COD-AB polygons $\rightarrow$ join rules
    $\rightarrow$ unit points $\rightarrow$ event centroid $\rightarrow$ GraphDB $\rightarrow$ map}}
  \caption{Method overview.}
  \Description{Pipeline from administrative boundary data to event points in the knowledge graph.}
  \label{fig:overview}
\end{figure}

\subsection{Data Sources}

Table~\ref{tab:sources} lists the data used. PSGC codes are 10-digit hierarchical identifiers
(region, province, city/municipality, barangay). The OCHA Common Operational Dataset of
administrative boundaries (COD-AB)~\cite{ocha2025codab} identifies features by \emph{pcodes},
the PSGC digits prefixed with \texttt{PH} and without trailing zeros (e.g.\ \texttt{PH0102801}
is PSGC \texttt{0102801000}). The two datasets have different vintages: the boundaries are valid
as of February 2025, while the PSGC release~\cite{psa2026psgc} is from the third quarter of
2026. Administrative changes between the two dates are the main difficulty in joining them.

\begin{table*}[t]
  \caption{Data sources.}
  \label{tab:sources}
  \small
  \begin{tabular}{p{4.2cm}p{4.4cm}p{3.0cm}p{3.6cm}}
    \toprule
    Source & Content used & Version & Role \\
    \midrule
    PSA PSGC publication datafile & codes, names, levels and hierarchy of all units & 3rd Quarter 2026 & location backbone (43,766 units) \\
    OCHA/HDX COD-AB Philippines (sources PSA, NAMRIA) & province (88), city/municipality (1,642) and barangay polygons & v03, valid 2025-02-13 & geometry for centroids; CC BY-IGO 3.0 \\
    DSWD DROMIC situation reports & per-municipality affected and displaced families and persons, damaged houses, assistance & reports published 2026 & event records and centroid weights \\
    \bottomrule
  \end{tabular}
\end{table*}

\subsection{Administrative-Unit Coordinates}

\paragraph{Centroid computation.} Each COD-AB polygon is projected from WGS~84 (EPSG:4326) to
UTM zone 51N (EPSG:32651), so that areas and centroids are computed in metres. For a multipart
polygon only the largest part is used, so that the point falls on the main landmass rather than
in the sea between islands. The area-weighted centroid of that part is taken; if it falls
outside the part (concave shapes), Shapely's \texttt{representative\_point()}, which is
guaranteed to lie inside, is used instead. The point is projected back to WGS~84 and rounded to
six decimals. Compared with the centre point COD-AB publishes for each feature, the median
distance over 1,730 provinces, cities and municipalities is 1.0~km and the 95th percentile
5.5~km; the largest differences come from multipart units such as Palawan (313~km, pulled by the
distant Kalayaan Islands), where COD-AB centres the whole shape.

\paragraph{Joining boundaries to PSGC units.} Each COD-AB feature is matched to a PSGC unit by a
sequence of rules (Table~\ref{tab:join}). (1)~\emph{Code}: the converted pcode equals a PSGC code
\emph{and} the names agree (token-set similarity $\geq 85$). (2)~\emph{Region re-code}: units
moved to a new region after the boundary vintage (Negros Island Region, created 2024; Sulu) are
re-coded first. (3)~\emph{Name}: a unique normalised-name match within the region. Units still
without a point are filled from related units: (4)~\emph{parent} (the 14 sub-municipalities of
Manila take the city's point), (5)~\emph{barangay dissolve} (eight municipalities created in 2024
in the BARMM Special Geographic Area are rebuilt from their barangays' polygons), and
(6)~\emph{member dissolve} (regions are rebuilt from their current members, since the COD-AB
region layer predates the Negros Island Region).

\begin{table}[t]
  \caption{Units joined per rule.}
  \label{tab:join}
  \small
  \begin{tabular}{lr}
    \toprule
    Rule & Units \\
    \midrule
    Code (with name agreement) & 1,562 \\
    Region re-code & 85 \\
    Name within region & 69 \\
    Member dissolve (regions) & 18 \\
    Parent (Manila districts) & 14 \\
    Barangay dissolve & 8 \\
    \midrule
    Total & 1,756 of 1,756 \\
    \bottomrule
  \end{tabular}
\end{table}

The name-agreement requirement in rule~1 is essential. In Maguindanao del Norte and del Sur,
renumbered after the boundary vintage, 16 COD-AB codes now identify a \emph{different} PSGC
municipality (e.g.\ COD-AB \texttt{PH1908822} is Pagagawan while PSGC \texttt{1908822000} is
South Upi). A code-only join would have placed those towns' points in neighbouring towns without
raising any error.

\paragraph{Representation.} Coordinates follow GeoSPARQL~1.1~\cite{ogc2024geosparql}: each unit
has a geometry node with a WKT point in longitude--latitude order, plus plain
\texttt{:latitude}/\texttt{:longitude} values. \texttt{:geoPrecision} records the level of the
polygon the point came from, and each geometry is linked with PROV-O~\cite{w3c2013provo} to a
dataset node recording the COD-AB version, validity date and licence:
{\scriptsize
\begin{verbatim}
:0500503000 geo:hasGeometry :0500503000/geometry .
:0500503000/geometry a geo:Geometry ;
  geo:asWKT "POINT (123.688861 13.114813)"^^geo:wktLiteral ;
  :geoPrecision "Municipality" ;
  prov:wasDerivedFrom :dataset/cod-ab-phl-v03 .
\end{verbatim}
}

\subsection{Event Centroids}

A DROMIC report tabulates the impact of an event per affected municipality. The centroid uses
the first available of: (a)~the affected-population rows, which carry counts; (b)~the location
named in the report title; (c)~the housing-damage rows. For an event with $n$ distinct locations
with points $(\varphi_i, \lambda_i)$ and weights $w_i$:

\begin{enumerate}
\item \textbf{Point lookup.} A location without its own point (a barangay) takes the point of its
  nearest ancestor in the PSGC hierarchy.
\item \textbf{Weights.} One count field is chosen for the whole event: the first of affected
  persons, affected families, displaced persons and displaced families that is positive for
  every location. Fields are never mixed; if none covers every location, $w_i = 1$. Rows for the
  same location are merged by summing.
\item \textbf{Geographic mean.} Points are averaged as unit vectors:
  \begin{gather*}
    x_i = \cos\varphi_i\cos\lambda_i,\quad y_i = \cos\varphi_i\sin\lambda_i,\quad z_i = \sin\varphi_i,\\
    X = \textstyle\sum w_i x_i,\quad Y = \textstyle\sum w_i y_i,\quad Z = \textstyle\sum w_i z_i,\\
    \bar\varphi = \operatorname{atan2}\bigl(Z, \sqrt{X^2+Y^2}\bigr),\qquad \bar\lambda = \operatorname{atan2}(Y, X).
  \end{gather*}
  A plain average of latitudes and longitudes differs from this by about 1~m for neighbouring
  towns, about 16~m across a region and up to 1.7~km for a country-wide event.
\item \textbf{Precision.} The coarsest precision among the inputs (an event with one
  province-only location is labelled Province).
\item \textbf{Spread.} The largest haversine~\cite{sinnott1984haversine} distance from the event
  point to any input point.
\end{enumerate}

The unweighted mean is always computed as well; the distance between the two (\emph{weighting
shift}) is reported. The event geometry carries \texttt{:centroidMethod},
\texttt{:locationCount} and \texttt{:spreadKm} in addition to the WKT point and precision.

\paragraph{Results on real reports.} Table~\ref{tab:events} shows three reports of different
shapes processed end to end. The El Ni\~no point lies in central Panay, pulled toward Tapaz, the
most affected town (22,795 persons); the Mayon point lies about 17~km south-west of the crater,
between the two most affected towns, Guinobatan (102,147) and Ligao (80,926), with all 13 towns
resolved. The fire's first run placed it in the City of Manila, 12.0~km away: the location matcher
tried the province column (``Metro Manila'') as a city before the municipality column. The defect
was visible only because the report title and the table named different places.

\begin{table*}[t]
  \caption{Event centroids for three DROMIC reports.}
  \label{tab:events}
  \footnotesize
  \setlength{\tabcolsep}{3pt}
  \begin{tabular}{p{3.4cm}p{3.0cm}lllrr}
    \toprule
    Event (report) & Locations & Point (lat, lon) & Method & Precision & Spread & Shift \\
    \midrule
    El Ni\~no, Western Visayas (Rep.~1) & 4 municipalities in Capiz, Iloilo & 11.174777, 122.454205 & weighted, persons & Municipality & 33.0 km & 8.8 km \\
    Mayon Volcano (Rep.~111) & 13 Albay municipalities and cities & 13.169899, 123.552245 & weighted, persons & Municipality & 38.1 km & 7.4 km \\
    Fire, Brgy.\ E.~Rodriguez (Rep.~4) & Quezon City & 14.680492, 121.060467 & single location & City & 0 & 0 \\
    \bottomrule
  \end{tabular}
\end{table*}

\subsection{Automated Collection}

DROMIC publishes each event as a WordPress post that accumulates a new report file at each
reporting cycle (typically 6~AM and 6~PM); a long-running event such as Mayon collects more than a
hundred files in one post. The original Selenium collector downloaded only the first file of a
post and never revisited it, so events stayed frozen at their first report.

The new collector reads the site's category RSS feed. Each item carries the post title (with
incident date and ``[Terminated]'' status), publication time, categories naming the disaster type,
and the full post HTML, from which every report file is extracted, newest first. Change detection
is per file, using a manifest keyed by file URL. The feed is polled hourly, the update period it
declares; a conditional request (ETag/If-Modified-Since) makes an unchanged poll a single
HTTP~304 response. Each poll that finds new files runs the rest of the pipeline incrementally
(Figure~\ref{fig:update}): fetch, parse the new PDFs only, \emph{supersede} older reports of the
same event (event identifiers are minted from event name and start date, so only the newest report
is kept active), map and validate with SHACL~\cite{w3c2017shacl} while \emph{quarantining} events
that fail instead of blocking the whole update, and atomically replace the DROMIC named graph in
GraphDB. Two safeguards check that a download is a real document and that a place name from the
post title appears in the file name. Every run is logged, including the delay between a report's
publication and its appearance in the graph.

\begin{figure}[t]
  \centering
  \fbox{\parbox[c][3.0cm][c]{0.9\columnwidth}{\centering\small
    [Figure F12 placeholder] hourly RSS poll (304 $\rightarrow$ stop) $\rightarrow$ fetch
    $\rightarrow$ parse $\rightarrow$ supersede $\rightarrow$ map + SHACL (quarantine)
    $\rightarrow$ atomic graph replace $\rightarrow$ run log}}
  \caption{Automated update chain.}
  \Description{Stages of one automated update run.}
  \label{fig:update}
\end{figure}

\subsection{Evaluation Design}

Engineering correctness is tested with automated checks (Section~\ref{sec:testing}). Accuracy is
measured against a hand-built reference prepared by partners from the Department of Geography
with a written protocol that mirrors each pipeline rule. Comparing stage by stage attributes every
error to one stage (Table~\ref{tab:eval}). The centroid is compared twice: with the annotator's
hand-computed point, and with the point the pipeline's own formula gives from the
\emph{reference} rows, which isolates upstream (parsing and matching) errors from the centroid
method. A pilot uses 10 reports stratified by disaster type and table size.

\begin{table}[t]
  \caption{Accuracy measures per stage.}
  \label{tab:eval}
  \small
  \begin{tabular}{p{2.1cm}p{5.6cm}}
    \toprule
    Stage & Measure \\
    \midrule
    Collection & share of reports fetched with the correct file \\
    Metadata & exact match of event name, start date, report number \\
    Locations & recall and precision of PSGC codes (municipality level) \\
    Figures & exact-match rate, median relative error per count field \\
    Centroid & km to hand point; km to formula on reference rows \\
    \bottomrule
  \end{tabular}
\end{table}

\section{Implementation}
\label{sec:impl}

\subsection{Architecture and Tools}

The ETL pipeline is written in Python~3.14 and runs as stages that each read the previous stage's
files: fetch, parse (PDF to tables, Docling), transform (typed records with location, organisation
and disaster-type resolution), map (RDF via rdflib), validate (pySHACL) and load (GraphDB). This
work adds a geodata stage (GeoPandas with pyogrio, Shapely, pyproj, RapidFuzz) that runs once per
PSGC or boundary release, an event-centroid step inside mapping, the RSS collector and an update
orchestrator with hourly scheduling.

\subsection{Validation and Loading}

The geodata stage writes to a staged file that is promoted only if it conforms to the SHACL
shapes, so a failed run never replaces good data. The shapes require one WKT point per geometry,
a coordinate pair within a Philippine bounding box and a precision from a fixed list; the
event-geometry shape constrains the WKT numerically, so a swapped latitude/longitude always fails.

The graph is stored in Ontotext GraphDB~\cite{ontotext2025graphdb} with the OWL~2~RL ruleset.
Each source folder becomes a named graph, replaced in one RDF4J transaction so that a failed load
leaves the previous version in place. A separate development repository with the production
configuration keeps new data from overwriting the existing graph. The first full load (372k
explicit triples) took 5~min~10~s, dominated by inference over the PSGC graph. GeoSPARQL
functions work without extra configuration: \texttt{geof:distance} places the Mayon event point
17~km from the crater.

\subsection{Defects Found and Fixed}

Several silent defects were found during implementation, illustrating the risks of geographic
enrichment:
\begin{itemize}
\item \textbf{Code collisions after renumbering}: 16 Maguindanao towns would have been misplaced
  by a code-only join.
\item \textbf{Stale region polygons}: the boundary dataset's region layer put Region~VII's point
  on Negros Island, now another region.
\item \textbf{Metro Manila resolved to Manila}: every NCR row written ``X, Metro Manila, NCR'' was
  placed in the City of Manila. A first fix was rejected after an old-versus-new comparison on 34
  probe strings showed it broke barangay-below-city rows.
\item \textbf{Non-deterministic output}: precision labels and table-row identifiers changed
  between identical runs (tie order; unordered join), which matters once the pipeline re-runs
  automatically.
\item \textbf{Events frozen at their first report}: the original collector tracked progress per
  post and took only its first file.
\item \textbf{Same-day reports mixed}: the 6~AM and 6~PM reports of one day received the same
  folder name, so the second report's tables were written beside the first's.
\item \textbf{Report numbers never extracted}: the parser expected ``Report No.~111'' while
  DROMIC titles read ``Report \#111'', so every report was numbered 0.
\item \textbf{Events dropped on the next run}: a failed assistance table listed the whole event
  for re-parsing, so with hourly rebuilds an event appeared and then disappeared.
\end{itemize}

\section{Testing}
\label{sec:testing}

Testing follows the pipeline's stages in four levels (Table~\ref{tab:tests}). Unit tests use
synthetic inputs with known answers: for a 3:1 weighting the event point must lie one quarter of
the way between two towns, and the precision label must not depend on input order. The collector
is tested without network access against a fake server replaying a saved copy of the live feed,
covering an unchanged poll (HTTP~304), a failed download that must be retried, and a server
returning a web page instead of a PDF. The pipeline level contains regression probes for every
location-matching defect found, including those a rejected fix had broken, and checks that the
sample reports produce exactly the expected triples and event points.

\begin{table}[t]
  \caption{Test levels (all passing, 25 Nov 2026).}
  \label{tab:tests}
  \small
  \begin{tabular}{p{1.6cm}p{4.2cm}rr}
    \toprule
    Level & Scope & Tests & Time \\
    \midrule
    Unit & centroid arithmetic, feed parsing, change detection, supersede, accuracy measures & 68 & 2 s \\
    Pipeline & matcher probes, real-data centroids, SHACL, determinism, quarantine & 20 & 7 min \\
    Geodata & join and coordinate regression (7 checks) & 1 & 9 min \\
    Deployment & loaded repository; live feed & 2 & 10 s \\
    \bottomrule
  \end{tabular}
\end{table}

Writing the tests exposed two further errors: the accuracy script read a count printed
``1,053\textsuperscript{1}'' (footnote mark) as 10,531, because the regular-expression digit class
also matches superscript digits; and an earlier claim that a plain coordinate average agrees with
the spherical mean to under a metre held only for small events.

The accuracy harness was itself tested before use: a reference generated from the pipeline's own
output scores 100\% on every measure, and a reference with one altered count and one altered
municipality produces the expected drop in figure accuracy and location recall, listing both
errors. The Geography pilot on 10 reports is prepared (protocol, templates and reference points)
and its results will be reported in Section~\ref{sec:results}.

\section{Results and Discussion}
\label{sec:results}

% Preliminary (25 Nov 2026): pipeline output on the 10 pilot reports; the Geography
% reference and the GraphDB load of these events are pending.

All 10 pilot reports were parsed, mapped and validated by the automated update in one run:
SHACL validation passed (3,271 triples), no event was quarantined, and every affected location
resolved to a PSGC unit (Table~\ref{tab:pilot}). All five fires were placed in the correct city or
municipality, including both Metro Manila cases, which an earlier version had placed in the City of
Manila, and a city whose report spelling differs from the PSGC (Ozamis/Ozamiz).

\begin{table}[t]
  \caption{Event centroids for the 10 pilot reports (preliminary).}
  \label{tab:pilot}
  \footnotesize
  \setlength{\tabcolsep}{3pt}
  \begin{tabular}{p{2.6cm}rp{1.5cm}rr}
    \toprule
    Event & Loc. & Method & Spread & Shift \\
    \midrule
    Dengue, Bicol (R10) & 106 & unweighted & 172.8 km & 0 \\
    El Ni\~no, national (R42) & 122 & persons & 824.8 km & 156.4 km \\
    El Ni\~no, W. Visayas (R1) & 4 & persons & 33.0 km & 8.8 km \\
    Thunderstorms, W. Visayas & 2 & persons & 30.7 km & 14.2 km \\
    Mayon Volcano (R111) & 13 & persons & 38.1 km & 7.4 km \\
    Five fires & 1 each & single & 0 & 0 \\
    \bottomrule
  \end{tabular}
\end{table}

A single point for a nationwide event is not meaningful. The national El Ni\~no point lands in
western Camarines Norte, a town with no special relation to a drought affecting 122 municipalities
from Cagayan to Mindanao; its spread (825~km) and weighting shift (156~km) flag the case. For such
events a map should show the affected municipalities rather than, or in addition to, the event
point. In the same report, 86 of 122 rows give identical counts for affected families and persons;
this is in the source document, not a parsing error. The dengue report, in contrast, exposed
genuine parsing errors: shifted rows, and a merged cell that turned ``3,351 3,494'' into 33,513,494
affected persons. A plausibility rule now rejects person counts above 20 per family. The
Geography pilot (Section~\ref{sec:method}) will measure how often such errors occur.

\section{Recommendations}

lorem ipsum  lorem ipsum  lorem ipsum  lorem ipsum  lorem ipsum  lorem ipsum  lorem ipsum  lorem ipsum  lorem ipsum  lorem ipsum  lorem ipsum  lorem ipsum  lorem ipsum  lorem ipsum  lorem ipsum  lorem ipsum  lorem ipsum  lorem ipsum  lorem ipsum  lorem ipsum  lorem ipsum  lorem ipsum  lorem ipsum  lorem ipsum  lorem ipsum  lorem ipsum  lorem ipsum  lorem ipsum  lorem ipsum  lorem ipsum  lorem ipsum  lorem ipsum  lorem ipsum  lorem ipsum  lorem ipsum  lorem ipsum  lorem ipsum  lorem ipsum  lorem ipsum  lorem ipsum  lorem ipsum  lorem ipsum  lorem ipsum  lorem ipsum  lorem ipsum  lorem ipsum  lorem ipsum  lorem ipsum  lorem ipsum  lorem ipsum  lorem ipsum  

\section{Conclusion}

lorem ipsum  lorem ipsum  lorem ipsum  lorem ipsum  lorem ipsum  lorem ipsum  lorem ipsum  lorem ipsum  lorem ipsum  lorem ipsum  lorem ipsum  lorem ipsum  lorem ipsum  lorem ipsum  lorem ipsum  lorem ipsum  lorem ipsum  lorem ipsum  lorem ipsum  lorem ipsum  lorem ipsum  lorem ipsum  lorem ipsum  lorem ipsum  lorem ipsum  lorem ipsum  lorem ipsum  lorem ipsum  lorem ipsum  lorem ipsum  lorem ipsum  lorem ipsum  lorem ipsum  lorem ipsum  lorem ipsum  lorem ipsum  lorem ipsum  lorem ipsum  lorem ipsum  lorem ipsum  lorem ipsum  lorem ipsum  lorem ipsum  lorem ipsum  lorem ipsum  lorem ipsum  lorem ipsum  lorem ipsum  lorem ipsum  lorem ipsum  lorem ipsum  

\section{Future Works}

\begin{itemize}
\item Barangay-level precision: COD-AB already provides barangay polygons, and DROMIC fire reports
  name a single barangay.
\item Population-weighted unit centroids, using the 2024 population already in the PSGC graph.
\item Per-report event snapshots, since DROMIC events evolve over many reports.
\item Unit polygons in the graph, enabling GeoSPARQL containment queries.
\item A larger accuracy study with two annotators per report and inter-annotator agreement.
\item Incremental mapping (one RDF file per event) instead of rebuilding a year's file per update.
\item A title-versus-table location consistency check as an automatic validation rule.
\item Conversion of Word reports without Microsoft Word (e.g.\ LibreOffice), so the pipeline can
  run on a Linux server.
\end{itemize}

\section{Acknowledgments}

Administrative boundaries: OCHA/HDX COD-AB Philippines (v03, valid 2025-02-13; sources PSA and
NAMRIA), licensed CC BY-IGO 3.0; centroids derived by the authors. Administrative codes: PSA
Philippine Standard Geographic Code. Disaster reports: DSWD DROMIC.
% TODO (authors): advisers, Department of Geography partners, SakunaGraPH team.

\section{Citations and Bibliographies}

lorem ipsum  lorem ipsum  lorem ipsum  lorem ipsum  lorem ipsum  lorem ipsum  lorem ipsum  lorem ipsum  lorem ipsum  lorem ipsum  lorem ipsum  lorem ipsum  lorem ipsum  lorem ipsum  lorem ipsum  lorem ipsum  lorem ipsum  lorem ipsum  lorem ipsum  lorem ipsum  lorem ipsum  lorem ipsum  lorem ipsum  lorem ipsum  lorem ipsum  lorem ipsum  lorem ipsum  lorem ipsum  lorem ipsum  lorem ipsum  lorem ipsum  lorem ipsum  lorem ipsum  lorem ipsum  lorem ipsum  lorem ipsum  lorem ipsum  lorem ipsum  lorem ipsum  lorem ipsum  lorem ipsum  lorem ipsum  lorem ipsum  lorem ipsum  lorem ipsum  lorem ipsum  lorem ipsum  lorem ipsum  lorem ipsum  lorem ipsum  lorem ipsum  

\section{Appendices}

lorem ipsum  lorem ipsum  lorem ipsum  lorem ipsum  lorem ipsum  lorem ipsum  lorem ipsum  lorem ipsum  lorem ipsum  lorem ipsum  lorem ipsum  lorem ipsum  lorem ipsum  lorem ipsum  lorem ipsum  lorem ipsum  lorem ipsum  lorem ipsum  lorem ipsum  lorem ipsum  lorem ipsum  lorem ipsum  lorem ipsum  lorem ipsum  lorem ipsum  lorem ipsum  lorem ipsum  lorem ipsum  lorem ipsum  lorem ipsum  lorem ipsum  lorem ipsum  lorem ipsum  lorem ipsum  lorem ipsum  lorem ipsum  lorem ipsum  lorem ipsum  lorem ipsum  lorem ipsum  lorem ipsum  lorem ipsum  lorem ipsum  lorem ipsum  lorem ipsum  lorem ipsum  lorem ipsum  lorem ipsum  lorem ipsum  lorem ipsum  lorem ipsum  

\section*{Ethics and Privacy Statement}

lorem ipsum  lorem ipsum  lorem ipsum  lorem ipsum  lorem ipsum  lorem ipsum  lorem ipsum  lorem ipsum  lorem ipsum  lorem ipsum  lorem ipsum  lorem ipsum  lorem ipsum  lorem ipsum  lorem ipsum  lorem ipsum  lorem ipsum  lorem ipsum  lorem ipsum  lorem ipsum  lorem ipsum  lorem ipsum  lorem ipsum  lorem ipsum  lorem ipsum  lorem ipsum  lorem ipsum  lorem ipsum  lorem ipsum  lorem ipsum  lorem ipsum  lorem ipsum  lorem ipsum  lorem ipsum  lorem ipsum  lorem ipsum  lorem ipsum  lorem ipsum  lorem ipsum  lorem ipsum  lorem ipsum  lorem ipsum  lorem ipsum  lorem ipsum  lorem ipsum  lorem ipsum  lorem ipsum  lorem ipsum  lorem ipsum  lorem ipsum  lorem ipsum  

%%
%% The next two lines define the bibliography style to be used, and
%% the bibliography file.
\bibliographystyle{ACM-Reference-Format}
\bibliography{references}


%%
%% If your work has an appendix, this is the place to put it.
\appendix

\section{Research Methods}

\subsection{Part One}

\subsection{Part Two}

\section{Online Resources}

\end{document}
\endinput
%%
%% End of file `acmtog.tex'.
